% ============================================================================
% Fichier  : partie6/P06_C26_trilemme_CRO.tex
% Titre    : Le Trilemme CRO -- Formalisation et Implications
% Auteur   : MINKA MI NGUIDJOÏ Thierry Emmanuel
% ============================================================================

\chapter{Le Trilemme CRO --- Formalisation et Implications}
\label{chap:CRO}

\epigraph{%
    « Security is always a trade-off between confidentiality, integrity,
    and availability. The perfect balance is the holy grail of our
    field. »%
}{--- Bruce Schneier}

\niveauChapitre{Recherche}{%
    Mathématiques niveau master recommandées (théorie de l'information~:
    entropie, information mutuelle). Ce chapitre formalise
    rétrospectivement l'outil utilisé intuitivement dans chaque
    \texttt{boxCRO} depuis le Prologue~: relire quelques exemples
    antérieurs avant d'aborder ce chapitre est conseillé.%
}

\begin{boxobjectifs}
\begin{itemize}
    \item Maîtriser les fondements de théorie de l'information
          nécessaires~: entropie de Shannon, information mutuelle,
          distance en variation totale, inégalité de Pinsker.
    \item Comprendre la formalisation rigoureuse du cadre CLO
          (\textit{Cryptographic Legal Opposability})~: définitions
          précises de $\mathsf{Conf}$, $\mathsf{Rel}$ et $\mathsf{Opp}$.
    \item Énoncer et comprendre la preuve du \textbf{théorème
          d'impossibilité} du Trilemme CRO conditionnel, établi
          par inégalités de Pinsker et de Fano.
    \item Analyser les trois études de cas documentées (Helios,
          zk-SNARKs, ECDSA) et comprendre pourquoi elles illustrent
          des points distincts de la surface de compromis.
    \item Relier la borne théorique aux limites cognitives humaines
          (capacité d'extraction institutionnelle $\Delta$) et à la
          littérature en sciences cognitives.
    \item Appliquer ce cadre théorique, désormais établi, à l'analyse
          de systèmes forensiques concrets rencontrés dans ce cours.
\end{itemize}
\end{boxobjectifs}

\bigskip

% ============================================================================
\section*{Un résultat établi, pas une conjecture}
% ============================================================================

Depuis le Prologue de ce cours, chaque chapitre a mobilisé une
\texttt{boxCRO} pour analyser une décision investigative~: fallait-il
collecter l'intégralité du téléphone de MBONGO ou seulement les
données pertinentes (Ch.~\ref{chap:acq})~? Comment arbitrer entre
conservation urgente d'une preuve cloud et respect de la vie privée
(Ch.~\ref{chap:norm-glob})~? Pourquoi la chaîne de custody
maximise-t-elle l'opposabilité au prix d'une contrainte sur la
confidentialité (Ch.~\ref{chap:custody})~?

Ce chapitre formalise l'outil que vous avez utilisé tout au long
du cours~: le \textbf{Trilemme CRO}. Il ne s'agit plus, à la date
de rédaction de ce cours, d'une proposition spéculative~: le cadre
a été \textbf{publié dans IEEE Access}, revue à comité de lecture
indexée (facteur d'impact~$\approx 3{,}6$, classement Clarivate
JCR), en deux articles complémentaires~:

\begin{itemize}
    \item \textbf{CLO Framework}~\cite{minka2026clo} établit les
          définitions formelles de $\mathsf{Conf}$, $\mathsf{Rel}$
          et $\mathsf{Opp}$, prouve les bornes fondamentales
          (Théorème~\ref{thm:bounds} et Théorème~\ref{thm:bottleneck}
          de ce chapitre) et documente trois études de cas réelles~;
    \item \textbf{Conditional CRO Trilemma}~\cite{minka2026conditional}
          établit la \textbf{borne générale d'impossibilité} restée
          ouverte dans le premier article~--- le résultat central de
          ce chapitre (Théorème~\ref{thm:trilemma-general}).
\end{itemize}

Ce chapitre n'est donc pas un exposé de spéculation théorique~:
c'est la présentation pédagogique d'un résultat de recherche
publié et évalué par les pairs, replacé dans le contexte appliqué
de l'investigation numérique qui occupe ce cours.

\separateur

% ============================================================================
\section{Fondements de Théorie de l'Information}
\label{sec:fondements-ti}
% ============================================================================

La formalisation du Trilemme CRO repose entièrement sur la théorie
de l'information de Shannon~\cite{shannon1948}. Cette section pose
les quatre outils mathématiques nécessaires, avec un niveau de
détail suffisant pour suivre les démonstrations sans exiger un
cours de théorie de l'information complet.

\subsection{Entropie de Shannon}

\begin{boxdef}
\textbf{Entropie}\index{entropie de Shannon}~: pour une variable
aléatoire discrète $X$ de loi $p_X$, l'entropie de Shannon mesure
l'incertitude moyenne (en bits) associée à $X$~:
\begin{equation}
H(X) = -\sum_{x \in \mathcal{X}} p_X(x) \log_2 p_X(x)
\label{eq:entropie}
\end{equation}
L'\textbf{entropie conditionnelle} $H(X \mid Y)$ quantifie
l'incertitude résiduelle sur $X$ après observation de $Y$.
\end{boxdef}

\textbf{Intuition forensique}~: $H(V \mid C)$, l'entropie d'un
verdict $V$ (par exemple, une qualification judiciaire) sachant le
contexte institutionnel $C$ (les compétences du juge, le cadre
légal applicable), mesure le nombre de scénarios légaux distincts
qu'un décideur peut raisonnablement distinguer compte tenu de ce
qu'il sait déjà.

\subsection{Information mutuelle}

\begin{boxdef}
\textbf{Information mutuelle}~: mesure la quantité d'information
partagée entre deux variables aléatoires~:
\begin{equation}
I(X;Y) = H(X) - H(X \mid Y) = H(Y) - H(Y \mid X)
\label{eq:info-mutuelle}
\end{equation}
$I(X;Y) = 0$ signifie que $X$ et $Y$ sont indépendants~: observer
$Y$ ne renseigne en rien sur $X$. $I(X;Y) = H(X)$ signifie que $Y$
détermine entièrement $X$.
\end{boxdef}

C'est cette quantité, appliquée aux preuves cryptographiques $\Sigma$
et aux verdicts institutionnels $V$, qui formalise précisément ce
que ce cours a appelé intuitivement ``l'opposabilité'' depuis le
Ch.~\ref{chap:custody}.

\subsection{Distance en variation totale}

\begin{boxdef}
\textbf{Distance en variation totale}~: pour deux distributions de
probabilité $P, Q$ sur le même espace~:
\begin{equation}
\mathrm{TV}(P,Q) = \frac{1}{2}\sum_{x} |P(x) - Q(x)|
= \max_{A} |P(A) - Q(A)|
\label{eq:tv}
\end{equation}
$\mathrm{TV}(P,Q) = 0$ si et seulement si $P = Q$ (distributions
parfaitement indistinguables)~; $\mathrm{TV}(P,Q) \leq 1$ toujours.
\end{boxdef}

\textbf{Intuition forensique}~: si deux messages $m_0$ et $m_1$
produisent des preuves cryptographiques dont les distributions ont
une distance en variation totale nulle, alors aucun observateur,
quels que soient ses moyens, ne peut déterminer lequel des deux
messages a réellement été chiffré. C'est précisément la définition
de la confidentialité parfaite mobilisée dans le cadre CLO.

\subsection{Inégalité de Pinsker}

\begin{boxjuris}
\textbf{Lemme (Inégalité de Pinsker)}~: pour deux distributions
$P, Q$~:
\begin{equation}
\mathrm{TV}(P,Q) \leq \sqrt{\tfrac{1}{2}\,D_{\mathrm{KL}}(P \| Q)}
\label{eq:pinsker}
\end{equation}
où $D_{\mathrm{KL}}$ est la divergence de Kullback-Leibler.
Cette inégalité \textbf{relie une mesure de distinguabilité
statistique} (variation totale) \textbf{à une mesure
d'information} (divergence, liée à l'information mutuelle).
C'est l'outil central qui permet, dans~\cite{minka2026conditional},
de traduire une exigence de fiabilité (haute information mutuelle
entre message et preuve) en une borne supérieure sur la
confidentialité atteignable.
\end{boxjuris}

\subsection{Inégalité de Fano}

\begin{boxjuris}
\textbf{Lemme (Inégalité de Fano)}~: pour un estimateur $\hat M$
d'une variable $M$ à partir d'une observation, avec probabilité
d'erreur $p_e = \Pr[\hat M \neq M]$~:
\begin{equation}
H(M \mid \hat M) \leq h(p_e) + p_e \log_2(|\mathcal{M}|-1)
\label{eq:fano}
\end{equation}
où $h(p_e)$ est l'entropie binaire. Cette inégalité formalise une
idée intuitive~: \textbf{si l'on peut deviner $M$ avec une faible
erreur, alors $M$ ne peut pas être trop incertain compte tenu de
l'observation}. Appliquée au contexte forensique, elle force la
fiabilité ($\mathsf{Rel}$ élevé, faible erreur de vérification) à
impliquer une information mutuelle minimale entre message et
preuve~--- l'un des deux verrous de la démonstration du théorème
principal de ce chapitre.
\end{boxjuris}

% ============================================================================
\section{Le Cadre CLO~: Définitions Formelles}
\label{sec:cadre-clo}
% ============================================================================

\subsection{Le protocole CLO}

\begin{boxdef}
\textbf{Protocole CLO} (\textit{Cryptographic Legal Opposability})
\cite{minka2026clo}~: un protocole $\mathcal{P}$ consiste en~:
\begin{itemize}
    \item des espaces~: messages $\mathcal{M}$, preuves $\Sigma$,
          contextes $\mathcal{C}$, verdicts $\mathcal{V}$~;
    \item des algorithmes~: $\mathsf{Setup}$, $\mathsf{Prove}$,
          $\mathsf{Verify}$, et surtout
          $\mathsf{Interpret}(\sigma, c) \to v$~: \textbf{l'algorithme
          d'interprétation institutionnelle} qui transforme une
          preuve cryptographique $\sigma$, dans un contexte $c$,
          en un verdict exploitable $v$.
\end{itemize}
\end{boxdef}

L'apport conceptuel central de ce cadre est de formaliser
explicitement ce que la cryptographie classique ignore~:
$\mathsf{Verify}$ garantit la correction mathématique~;
$\mathsf{Interpret}$ formalise la capacité d'un acteur institutionnel
(juge, auditeur, régulateur~--- ou, dans ce cours, l'OPJ et le
magistrat camerounais) à \textbf{extraire un sens exploitable} de
cette preuve. C'est cette seconde exigence, absente des définitions
de sécurité cryptographique standard, que ce cours a systématiquement
appelée ``opposabilité''.

\subsection{Les trois mesures}

\begin{boxdef}
\textbf{Confidentialité}~\cite{minka2026clo}~:
\begin{equation}
\mathsf{Conf}(\mathcal{P}) = 1 - \max_{m_0,m_1 \in \mathcal{M}}
\mathrm{TV}\bigl(P_{\Sigma \mid M=m_0},\, P_{\Sigma \mid M=m_1}\bigr)
\label{eq:conf}
\end{equation}
$\mathsf{Conf} = 1$~: les preuves sont parfaitement indistinguables
quel que soit le message (confidentialité parfaite, au sens de la
théorie de l'information). $\mathsf{Conf} = 0$~: les preuves
révèlent complètement le message.

\textbf{Fiabilité}~\cite{minka2026clo}~:
\begin{equation}
\mathsf{Rel}(\mathcal{P}) = \min_{m \in \mathcal{M}, c \in \mathcal{C}}
\Pr[\mathsf{Verify}(pk, \mathsf{Prove}(sk,m,c), c) = 1]
\label{eq:rel}
\end{equation}
$\mathsf{Rel} = 1$~: les preuves honnêtes sont toujours vérifiées
correctement (complétude parfaite).

\textbf{Opposabilité}~\cite{minka2026clo}~:
\begin{equation}
\mathsf{Opp}(\mathcal{P}, C) = \frac{I(V; \Sigma \mid C)}{H(V \mid C)}
\label{eq:opp}
\end{equation}
$\mathsf{Opp}$ mesure la \textbf{fraction de l'incertitude sur le
verdict que l'observation de la preuve permet effectivement de
lever}, étant donné le contexte institutionnel. $\mathsf{Opp} = 1$~:
la preuve détermine entièrement le verdict (opposabilité parfaite).
$\mathsf{Opp} = 0$~: la preuve n'apporte aucune information
exploitable au-delà du contexte déjà disponible.
\end{boxdef}

\begin{boxalert}
\textbf{Sur la nature mesurable de $\mathsf{Opp}$~: une innovation
du cadre CLO}

Contrairement à la confidentialité et à la fiabilité, déjà
formalisées de longue date en cryptographie, l'opposabilité au sens
de la définition~\eqref{eq:opp} \textbf{est une mesure originale}
introduite dans~\cite{minka2026clo}. C'est elle qui permet de
quantifier rigoureusement une intuition que ce cours a mobilisée
de façon qualitative depuis le Prologue~: une preuve peut être
mathématiquement irréprochable ($\mathsf{Rel}$ élevé) tout en étant
\textbf{institutionnellement inexploitable} si aucun décideur ne
peut en extraire un verdict actionnable~--- exactement le
``\textit{Invisible Authenticity Paradox}'' documenté dans
l'article fondateur (Section~\ref{sec:etudes-cas}).
\end{boxalert}

% ============================================================================
\section{Les Théorèmes Fondamentaux}
\label{sec:theoremes}
% ============================================================================

\subsection{Bornes de mesure et goulot d'étranglement contextuel}

\begin{boxjuris}
\textbf{Théorème~1 (Bornes de mesure)}~\cite{minka2026clo}~: pour
tout protocole CLO $\mathcal{P}$~:
\begin{equation}
0 \leq \mathsf{Conf}(\mathcal{P}) \leq 1, \quad
0 \leq \mathsf{Rel}(\mathcal{P}) \leq 1, \quad
0 \leq \mathsf{Opp}(\mathcal{P}, \mu) \leq 1
\label{thm:bounds}
\end{equation}
\end{boxjuris}

\begin{boxjuris}
\textbf{Théorème~2 (Goulot d'étranglement contextuel)}~%
\cite{minka2026clo}~: pour tout protocole CLO avec fonction
d'interprétation $\mathsf{Interpret}$ et distribution de contexte
$\mu$ sur $\mathcal{C}$~:
\begin{equation}
\mathsf{Opp}(\mathcal{P}, C) \leq \min\left\{1,\;
\frac{H(V \mid C)}{H(V \mid C)}\right\}
= \min\left\{1, \frac{I(V;\Sigma\mid C)}{H(V\mid C)}\right\}
\label{thm:bottleneck}
\end{equation}
\textbf{Interprétation}~: l'opposabilité est structurellement
plafonnée par la \textbf{capacité de connaissance de
l'interprète}~--- c'est-à-dire par ce que le contexte institutionnel
($H(C)$~: la formation du juge, la richesse du cadre légal,
l'expertise de l'auditeur) permet effectivement de distinguer.
Aucune sophistication cryptographique supplémentaire ne peut
compenser un goulot d'étranglement de compréhension institutionnelle.
\end{boxjuris}

\begin{boiteRemarque}{Une lecture directement applicable au Cameroun}
Le Théorème~2 explique, en termes rigoureux, une observation
qualitative déjà rencontrée au Ch.~\ref{chap:droit-cam}~: la
formation des magistrats camerounais au droit du numérique
constitue un facteur limitant pour l'opposabilité des preuves
numériques, indépendamment de la qualité technique de l'expertise
forensique produite. Investir dans $H(C)$ (la formation des juges,
des OPJ, des avocats) augmente le plafond théorique d'opposabilité
atteignable par \emph{tout} système de preuve, là où investir
uniquement dans la sophistication cryptographique ($\Sigma$) ne le
fait pas.
\end{boiteRemarque}

\subsection{Le Théorème d'Impossibilité Général}

L'article fondateur~\cite{minka2026clo} avait établi le Corollaire~2
(la tension qualitative entre les trois dimensions) mais laissait
ouverte, explicitement, la question d'une \textbf{borne générale
quantifiée}. C'est cette question que résout
l'article~\cite{minka2026conditional}.

\begin{boxjuris}
\textbf{Théorème~3 (Trilemme CRO Conditionnel, pour protocoles
réguliers)}~\cite{minka2026conditional}~: sous des hypothèses de
capacité d'extraction bornée, d'espace de verdicts borné et de
régularité du protocole, pour tout protocole de preuve $\mathcal{P}$
avec paramètre de sécurité $\lambda \geq 64$~:
\begin{equation}
\mathsf{Conf}(\mathcal{P}) \cdot \mathsf{Rel}(\mathcal{P}) \cdot
\mathsf{Opp}(\mathcal{P}) \;\leq\; \frac{\Delta}{\lambda} +
O(\lambda^{-2}) + \mathrm{negl}(\lambda)
\label{thm:trilemma-general}
\end{equation}
où $\Delta$ quantifie la \textbf{capacité d'extraction
institutionnelle} (le nombre de bits d'information décisionnelle
qu'un interprète humain peut effectivement traiter), et le terme
$O(\lambda^{-2})$ capture les contraintes cryptographiques d'ordre
supérieur, négligeables pour les paramètres de sécurité usuels
($\lambda \geq 128$).
\end{boxjuris}

\subsection{Esquisse de la démonstration}

La preuve complète, technique, occupe plusieurs pages
de~\cite{minka2026conditional}~; cette section en donne la structure
logique, suffisante pour en comprendre la portée sans en reproduire
tous les détails calculatoires.

\begin{enumerate}
    \item \textbf{La fiabilité force l'information mutuelle.}
          Par l'inégalité de Fano~\eqref{eq:fano}, si
          $\mathsf{Rel} \geq 1/2$ (toute preuve raisonnable doit
          au moins faire mieux qu'une pièce de monnaie), alors
          $I(M; \Sigma, C) \geq \mathsf{Rel} \cdot \lambda - 1$~:
          un système fiable nécessite que la preuve porte une
          quantité substantielle d'information sur le message.
    \item \textbf{L'opposabilité est bornée par l'extraction
          institutionnelle.} Si l'on modélise la capacité
          d'extraction humaine $K$ comme limitée à $\Delta$~bits
          (justifié en Section~\ref{sec:capacite-extraction}),
          alors $I(V;\Sigma\mid C) \lesssim I(M;\Sigma\mid C) + \Delta$.
    \item \textbf{L'espace des verdicts est borné.} Par le théorème
          de codage de source de Shannon, un espace de verdicts
          institutionnellement gérable impose
          $H(V \mid C) \geq \alpha \cdot H(C) - 1$ pour un paramètre
          $\alpha$ de granularité des verdicts.
    \item \textbf{La confidentialité est bornée par Pinsker.} Une
          forte information mutuelle $I(M;\Sigma\mid C)$ (requise
          pour la fiabilité, étape~1) implique, via
          l'inégalité~\eqref{eq:pinsker}, une borne supérieure sur
          $\mathsf{Conf}$~: plus la preuve est informative sur le
          message, moins elle peut rester confidentielle.
    \item \textbf{Combinaison et optimisation.} En combinant ces
          quatre bornes et en optimisant sur la variable
          $x = I(M;\Sigma\mid C)$ (un compromis interne entre
          fiabilité et confidentialité), on obtient le maximum
          atteignable du produit
          $\mathsf{Conf} \cdot \mathsf{Rel} \cdot \mathsf{Opp}$,
          qui se révèle être exactement~$\Delta/\lambda$ à des
          termes négligeables près.
\end{enumerate}

\begin{boxalert}
\textbf{Honnêteté scientifique sur la portée du résultat}

Comme tout résultat de recherche rigoureux, ce théorème s'applique
sous des \textbf{hypothèses explicites} (protocoles ``réguliers'',
capacité d'extraction additive et non multiplicative, espace de
verdicts borné)~--- discutées et justifiées
dans~\cite{minka2026conditional}, mais qui délimitent le domaine
de validité du résultat. Les auteurs eux-mêmes identifient des
formes alternatives de bornes (sous-linéaire, saturante) comme
pistes de recherche future. Cette transparence sur les limites
d'un résultat est elle-même une compétence forensique~:
documenter ce qu'une analyse établit \emph{et} ce qu'elle
n'établit pas (Ch.~\ref{chap:rapport}, distinction
constatation/analyse).
\end{boxalert}

% ============================================================================
\section{Capacité d'Extraction Institutionnelle~: le Paramètre $\Delta$}
\label{sec:capacite-extraction}
% ============================================================================

\subsection{Fondement en sciences cognitives}

Le paramètre $\Delta$ n'est pas un artifice mathématique sans
ancrage empirique~: il est directement relié aux travaux fondateurs
de Miller~\cite{miller1956} sur la capacité de la mémoire de travail
humaine, estimée à $7 \pm 2$ ``chunks'' d'information.

\begin{tableaucours}{p{3.5cm}p{3.0cm}p{4.5cm}}{%
    \tableauHeader{Composante} &
    \tableauHeader{Estimation} &
    \tableauHeader{Justification}
}
Capacité individuelle de base & $\approx 21$~bits &
    7~chunks~$\times$~3~bits par chunk (distinction entre~8~états
    possibles par chunk) \\
+ Documentation institutionnelle & $+5$ à $10$~bits &
    Procédures écrites, matériaux de référence
    (ex.~: guides de qualification, Ch.~\ref{chap:droit-cam}) \\
+ Collaboration & $+10$ à $20$~bits &
    Délibération multi-acteurs, mise en commun d'expertise
    (collège de juges, équipe d'enquête) \\
+ Spécialisation & $+10$ à $15$~bits &
    Formation spécifique au-delà de la cognition générale
    (expert agréé, Ch.~\ref{chap:droit-cam}) \\
\end{tableaucours}
\vspace{-0.8em}
\begin{center}{\small\itshape Tableau~26.1 --- Construction du paramètre
$\Delta$ pour une institution bien dotée}\end{center}

Cette construction additive conduit à une estimation
$\Delta \in [20, 60]$~bits pour des institutions bien dotées en
ressources, cohérente avec les analyses empiriques de contextes
réels présentées dans les deux articles.

\subsection{Estimations empiriques de $H(C)$ par contexte}

\begin{tableaucours}{p{3.5cm}p{2.5cm}p{5.0cm}}{%
    \tableauHeader{Contexte institutionnel} &
    \tableauHeader{$H(C)$ estimé} &
    \tableauHeader{Justification}
}
Tribunal de juridiction générale & $\approx 8$~bits &
    $\approx 256$~contextes évidentiaires distingués \\
Régulateur spécialisé & $\approx 11$~bits &
    $\approx 2048$~contextes (réglementation financière détaillée) \\
Expert technique de domaine & $\approx 20$~bits &
    $\approx 10^6$~contextes (connaissance fine du domaine) \\
Contexte électoral (Helios) & $30$ à $50$~bits &
    Registre des électeurs, structure du scrutin, règles du
    district \\
Conformité financière (zk-SNARKs) & $40$ à $70$~bits &
    Historique des transactions, réglementation, pistes d'audit \\
Signature contractuelle (ECDSA) & $20$ à $40$~bits &
    Parties, termes, juridiction~--- contexte volontairement
    restreint \\
\end{tableaucours}
\vspace{-0.8em}
\begin{center}{\small\itshape Tableau~26.2 --- Estimations empiriques de
$H(C)$ par type de contexte (d'après~\cite{minka2026clo,minka2026conditional})}\end{center}

\begin{boiteRemarque}{Le paradoxe de l'écart de capacité}
Une observation frappante des deux articles~: atteindre une
opposabilité significative ($\mathsf{Opp} \gtrsim 0{,}5$) pour des
paramètres de sécurité cryptographique usuels ($\lambda = 128$
à~$256$~bits) exigerait un contexte institutionnel
$H(C) \approx 64$ à~$128$~bits. Or les estimations empiriques de
contextes institutionnels réels (Tableau~26.2) plafonnent typiquement
à~$8$ à~$20$~bits~--- un \textbf{écart d'un facteur~5 à~10} entre
ce que la cryptographie moderne exige et ce que les institutions
peuvent effectivement traiter. C'est ce paradoxe que les auteurs
nomment l'\textit{Invisible Authenticity Paradox}.
\end{boiteRemarque}

% ============================================================================
\section{Trois Études de Cas Documentées}
\label{sec:etudes-cas}
% ============================================================================

L'article fondateur~\cite{minka2026clo} analyse trois systèmes
cryptographiques réellement déployés, révélant des positions
distinctes et cohérentes sur la surface de compromis CRO.

\subsection{Helios~: vérifiabilité électorale}

Le protocole de vote électronique Helios~\cite{adida2009helios}
combine chiffrement ElGamal et preuves à divulgation nulle pour
garantir l'intégrité du scrutin sans compromettre le secret du vote.

\begin{boxscenario}{Helios --- l'institutionnalisation impossible}
Lors de l'élection présidentielle de l'\textit{Université catholique
de Louvain} (mars~2009)~\cite{adida2009university}, les
``\textit{trustees}'' (fiduciaires) sélectionnés parmi des étudiants
et le personnel administratif~--- choisis pour leur représentativité
institutionnelle, pas pour leur expertise cryptographique~---
n'ont pas pu exécuter eux-mêmes les opérations de génération et de
déchiffrement des clés. Des experts cryptographiques externes ont
dû intervenir. Plus frappant encore~: l'\textit{International
Association for Cryptologic Research} (IACR), pourtant composée
de cryptographes professionnels, a dû \textbf{annuler entièrement}
son élection de novembre~2025 après qu'un fiduciaire a
irrémédiablement perdu sa portion de la clé de déchiffrement
distribuée~\cite{iacr2025}.
\end{boxscenario}

\textbf{Lecture CRO}~: Helios atteint une confidentialité et une
fiabilité élevées (le secret du vote et l'intégrité du dépouillement
sont mathématiquement garantis), mais une opposabilité
institutionnelle estimée à $\mathsf{Opp} \approx 0{,}50$~--- même
des cryptographes professionnels peinent à opérer le système de
façon autonome.

\subsection{zk-SNARKs~: conformité financière}

Les preuves à divulgation nulle succinctes (zk-SNARKs) permettent
à une institution financière de prouver sa conformité réglementaire
(lutte anti-blanchiment, filtrage des sanctions) via une preuve
de quelques centaines d'octets, sans révéler le détail des
transactions sous-jacentes.

\textbf{Lecture CRO}~: confidentialité quasi-parfaite des
transactions individuelles, mais lorsqu'un schéma suspect émerge
et qu'une investigation ciblée devient nécessaire (exactement le
type de situation rencontré au Ch.~\ref{chap:for-cloud} de ce
cours, analyse blockchain), l'opacité cryptographique qui protège
la confidentialité des clients \textbf{empêche précisément
l'investigation que la réglementation est censée permettre}.
L'article estime $\mathsf{Opp} \approx 0{,}67$ pour ce type de
système~--- meilleur qu'Helios, mais encore loin de l'optimum.

\subsection{ECDSA~: preuve juridique}

À l'opposé du spectre, les signatures numériques à courbes
elliptiques (ECDSA)~\cite{johnson2001ecdsa}, omniprésentes dans
les infrastructures à clé publique (Ch.~\ref{chap:leg-mond},
eIDAS), atteignent une \textbf{acceptation judiciaire quasi
universelle}~: les tribunaux comprennent l'affirmation sémantique
``la partie A a autorisé le document D'' sans formation
cryptographique, via des décennies de jurisprudence et
d'infrastructure PKI établie.

\textbf{Lecture CRO}~: $\mathsf{Opp} \approx 0{,}90$~--- la plus
élevée des trois systèmes étudiés~--- mais au prix d'une
confidentialité quasi nulle~: la vérification exige de révéler
publiquement les messages signés et les relations entre clés,
incompatible avec des contextes exigeant confidentialité
(contrats sensibles, documents judiciaires scellés).

\subsection{Synthèse}

\begin{tableaucours}{p{3.0cm}p{2.2cm}p{2.2cm}p{2.2cm}p{2.8cm}}{%
    \tableauHeader{Système} &
    \tableauHeader{$\mathsf{Conf}$} &
    \tableauHeader{$\mathsf{Rel}$} &
    \tableauHeader{$\mathsf{Opp}$} &
    \tableauHeader{Tension dominante}
}
Helios (vote) & Élevé & Élevé & $\approx 0{,}50$ &
    Confidentialité/Fiabilité $\to$ Opposabilité \\
zk-SNARK (finance) & Quasi-parfait & Élevé & $\approx 0{,}67$ &
    Confidentialité $\to$ Opposabilité (audit ciblé impossible) \\
ECDSA (juridique) & Quasi-nul & Élevé & $\approx 0{,}90$ &
    Opposabilité $\to$ Confidentialité \\
\end{tableaucours}
\vspace{-0.8em}
\begin{center}{\small\itshape Tableau~26.3 --- Synthèse des trois études
de cas (d'après~\cite{minka2026clo})~: les produits
$\mathsf{Conf}\cdot\mathsf{Rel}\cdot\mathsf{Opp}$ observés s'échelonnent
de~$0{,}27$ à~$0{,}57$ selon le système, cohérent avec la borne du
Théorème~3}\end{center}

% ============================================================================
\section{Implications pour la Conception de Systèmes Forensiques}
\label{sec:implications-conception}
% ============================================================================

\subsection{Du calcul optimal à l'ingénierie sous contrainte}

L'implication méthodologique majeure des deux
articles~\cite{minka2026clo,minka2026conditional} est un
\textbf{changement de paradigme de conception}~: puisqu'aucun
protocole ne peut maximiser simultanément les trois dimensions
(Théorème~3), la question pertinente n'est plus ``quelle est la
primitive optimale~?'' mais ``\textbf{quel arbitrage spécifique
ce contexte institutionnel exige-t-il, et quelle architecture
satisfait ce besoin~?}''

\begin{boxCRO}
\textbf{Analyse CRO --- L'architecture Gold-Iron-Clay
(\textit{Semantic One-Time Pad})}

L'article fondateur~\cite{minka2026clo} propose et prouve
optimale, au sein de sa classe de construction, une architecture
en trois couches~: \textbf{Gold} (confidentialité parfaite via
masque jetable), \textbf{Iron} (fiabilité via engagement
cryptographique liant), \textbf{Clay} (préservation contextuelle
pour l'opposabilité). Cette construction atteint
$\mathsf{Conf} \cdot \mathsf{Rel} \cdot \mathsf{Opp} =
H(C)/\lambda \cdot (1 - \mathrm{negl}(\lambda))$~--- l'optimum
exact permis par la borne du Théorème~2, démontrant qu'une
architecture multicouche peut effectivement \textbf{atteindre}
(et non simplement approcher) la limite théorique.

$\mathcal{R}$~: chaque couche apporte sa force dominante plutôt
que de chercher une primitive unique universelle~--- exactement
la logique de la chaîne de custody du Ch.~\ref{chap:custody}
(hash pour l'intégrité, signature pour l'authenticité, formulaire
papier pour la traçabilité).

$\mathcal{O}$~: la preuve d'optimalité, publiée et vérifiable,
confère à cette architecture une légitimité scientifique directement
opposable dans une discussion technique.

\cro{$\uparrow\uparrow$ couche Gold}{$\uparrow\uparrow$ couche Iron}{$\uparrow\uparrow$ couche Clay}
\end{boxCRO}

\subsection{Application à l'architecture forensique de ce cours}

\begin{tableaucours}{p{3.0cm}p{4.0cm}p{4.5cm}}{%
    \tableauHeader{Dimension prioritaire} &
    \tableauHeader{Mécanisme apportant cette force} &
    \tableauHeader{Exemple dans ce cours}
}
$\mathsf{Rel}$ (Fiabilité) & Hash cryptographique
    (SHA-256/512) & Formulaire de custody
    (Ch.~\ref{chap:custody})~: hash calculé avant et après
    acquisition \\
$\mathsf{Opp}$ (Opposabilité) & Signature qualifiée +
    horodatage + formation du décideur & eIDAS
    (Ch.~\ref{chap:leg-mond})~; investissement dans $H(C)$ via la
    formation des magistrats (Ch.~\ref{chap:droit-cam}) \\
$\mathsf{Conf}$ (Confidentialité) & Chiffrement +
    minimisation de la collecte & Extraction ADB ciblée plutôt
    qu'exhaustive (Ch.~\ref{chap:for-mob})~; proportionnalité
    (\loicam{2024/017}, art.~27) \\
\end{tableaucours}
\vspace{-0.8em}
\begin{center}{\small\itshape Tableau~26.4 --- Architecture en couches
appliquée au cadre forensique de ce cours}\end{center}

Cette table révèle une cohérence remarquable~: la pratique
forensique enseignée tout au long de ce cours~--- chaîne de custody,
distinction constatation/analyse, principe de proportionnalité~---
\textbf{constitue, sans que cela ait été formulé explicitement avant
ce chapitre, une instanciation pratique de l'architecture en couches
prouvée optimale par le cadre CLO}.

% ============================================================================
\section{Application Forensique Intégrative}
\label{sec:application-forensique}
% ============================================================================

\begin{boxscenario}{Relire l'affaire MVOM à travers le Trilemme CRO formalisé}
Reprenons l'Opération MVOM (Ch.~\ref{chap:for-sys},
\ref{chap:for-res}, \ref{chap:for-cloud}) à la lumière de ce
chapitre.

\textbf{Le hash \texttt{\$UsnJrnl}} (Ch.~\ref{chap:for-sys})~:
$\mathsf{Rel}$ proche de~1 (déterministe, vérifiable
indépendamment)~; $\mathsf{Conf}$ non pertinent (le hash ne protège
aucune confidentialité)~; $\mathsf{Opp}$ élevé \emph{si et seulement
si} le tribunal dispose de la formation technique pour interpréter
un horodatage de journal système~--- exactement la dépendance au
contexte $H(C)$ formalisée par le Théorème~2.

\textbf{Le chiffrement BitLocker du disque saisi}
(Ch.~\ref{chap:PIU})~: $\mathsf{Conf}$ élevé pour le suspect avant
saisie~; après saisie légale et obtention de la clé, le système
bascule vers un régime où $\mathsf{Rel}$ et $\mathsf{Opp}$
deviennent les dimensions pertinentes pour l'analyse forensique~---
illustrant que le positionnement CRO d'un mécanisme \textbf{dépend
du moment et du détenteur} de l'information, pas uniquement de
l'algorithme.

\textbf{La transaction Bitcoin tracée} (Ch.~\ref{chap:for-cloud})~:
$\mathsf{Conf}$ faible (registre public)~; $\mathsf{Rel}$ maximal
(immuabilité de la blockchain)~; $\mathsf{Opp}$ initialement faible
(pseudonymat) mais \textbf{augmentable} par l'apport de contexte
externe (identification KYC aux points d'échange)~--- une
illustration directe du Théorème~2~: l'opposabilité croît avec
$H(C)$, ici obtenu par corrélation avec des sources tierces plutôt
que par la primitive cryptographique elle-même.
\end{boxscenario}

\separateur

\begin{boxresume}
\textbf{Ce qu'il faut retenir de ce chapitre~:}
\begin{itemize}
    \item Le \textbf{Trilemme CRO} est désormais un résultat
          \textbf{publié et évalué par les pairs}
          (IEEE Access,~2026, IF~$\approx 3{,}6$)~: les
          Théorèmes~1, 2 et~3 de ce chapitre correspondent à des
          théorèmes prouvés dans~\cite{minka2026clo}
          et~\cite{minka2026conditional}, pas à des conjectures.

    \item \textbf{Outils de théorie de l'information mobilisés}~:
          entropie de Shannon, information mutuelle, distance en
          variation totale, inégalités de Pinsker et de Fano
          --- fondements classiques appliqués à un problème nouveau.

    \item \textbf{Définitions formelles}~: $\mathsf{Conf}$ (via TV
          distance), $\mathsf{Rel}$ (probabilité de vérification
          correcte), $\mathsf{Opp}$ (information mutuelle normalisée
          entre preuve et verdict)~--- cette dernière étant
          l'innovation centrale du cadre.

    \item \textbf{Théorème principal}~: $\mathsf{Conf} \cdot
          \mathsf{Rel} \cdot \mathsf{Opp} \leq \Delta/\lambda +
          O(\lambda^{-2})$, où $\Delta$ quantifie la capacité
          d'extraction institutionnelle, enracinée dans les travaux
          de Miller sur la mémoire de travail humaine.

    \item \textbf{Trois études de cas documentées}~: Helios
          ($\mathsf{Opp}\approx 0{,}50$), zk-SNARKs
          ($\approx 0{,}67$), ECDSA ($\approx 0{,}90$)~---
          chacune occupant un point distinct, cohérent avec la
          théorie, sur la surface de compromis.

    \item \textbf{Réponse pratique~: architecture en couches}
          (Gold-Iron-Clay), prouvée optimale dans sa classe de
          construction, dont ce cours a montré qu'elle correspond
          structurellement à la pratique forensique enseignée
          depuis le Ch.~\ref{chap:custody}.

    \item \textbf{Investir dans $H(C)$}~--- la formation des
          magistrats et des OPJ camerounais~--- élève le plafond
          d'opposabilité atteignable par \emph{tout} système de
          preuve, une implication directement actionnable pour
          la stratégie nationale évoquée au Ch.~\ref{chap:bench}.
\end{itemize}
\end{boxresume}

\bigskip

\begin{boxexo}[title={\ding{45}\quad Exercice~26.1 --- Maîtrise des outils de théorie de l'information \niveaubase}]
\begin{enumerate}[(a)]
    \item Calculez $H(X)$ pour une variable $X$ uniforme sur
          $4$~valeurs. Interprétez le résultat en bits.
    \item Si $I(X;Y) = H(X)$, que peut-on en conclure sur la
          relation entre $X$ et $Y$~?
    \item Expliquez, en une phrase, pourquoi l'inégalité de
          Pinsker~\eqref{eq:pinsker} permet de relier une exigence
          de fiabilité à une borne sur la confidentialité.
    \item Pour chacun des trois mécanismes suivants rencontrés
          dans ce cours, identifiez la dimension CRO structurellement
          privilégiée et celle structurellement sacrifiée, en
          justifiant par un argument technique~: le chiffrement
          BitLocker (Ch.~\ref{chap:PIU})~; une signature ECDSA
          (Ch.~\ref{chap:leg-mond})~; une preuve à divulgation
          nulle de connaissance.
\end{enumerate}
\end{boxexo}

\begin{boxexo}[title={\ding{45}\quad Exercice~26.2 --- Lecture critique du Théorème~3 \niveaumoyen}]
\begin{enumerate}[(a)]
    \item Le Théorème~3 s'applique sous l'hypothèse
          $\lambda \geq 64$. Pourquoi cette condition est-elle
          nécessaire (référez-vous au rôle du terme
          $O(\lambda^{-2})$)~?
    \item Le Tableau~26.2 estime $H(C) \approx 8$~bits pour un
          tribunal de juridiction générale. En utilisant le
          Théorème~2, et pour $\lambda = 128$, estimez l'opposabilité
          maximale théoriquement atteignable dans ce contexte si
          $H(V|C) \approx H(C)$. Commentez l'écart avec
          l'opposabilité ECDSA observée empiriquement
          ($\approx 0{,}90$)~--- que cela révèle-t-il sur les
          conditions d'application du théorème~?
    \item Le chapitre mentionne que la forme additive de $\Delta$
          (plutôt que multiplicative ou sous-linéaire) a été
          choisie et justifiée par les auteurs. Pourquoi une
          borne sous-linéaire serait-elle ``trop pessimiste'',
          selon les termes de l'article~?
\end{enumerate}
\end{boxexo}

\begin{boxexo}[title={\ding{45}\quad Exercice~26.3 --- Application intégrative à une affaire du cours \niveauexpert}]
Reprenez l'affaire DIBAMBA (Ch.~\ref{chap:cas-int})~: fraude MoMo
avec SIM~swap, communications chiffrées, et adresse Bitcoin tracée.
\begin{enumerate}[(a)]
    \item Pour chacun des trois éléments probatoires de cette
          affaire (logs MTN, communications chiffrées, transactions
          blockchain), estimez qualitativement (Faible/Modéré/Élevé,
          en justifiant) les trois mesures $\mathsf{Conf}$,
          $\mathsf{Rel}$, $\mathsf{Opp}$, en vous appuyant sur les
          définitions formelles de la Section~\ref{sec:cadre-clo}.
    \item En vous inspirant de l'architecture Gold-Iron-Clay
          (Section~\ref{sec:implications-conception}), construisez
          une architecture en couches pour cette affaire~: quel
          mécanisme apporte la force $\mathsf{Rel}$~? Lequel apporte
          $\mathsf{Opp}$~? Lequel gère $\mathsf{Conf}$~?
    \item Le Théorème~2 implique que l'opposabilité de cette
          affaire est plafonnée par $H(C)$ du tribunal saisi.
          Proposez deux actions concrètes, dans l'esprit du
          Tableau~26.1 (documentation, collaboration, spécialisation),
          qui augmenteraient $H(C)$ pour ce dossier spécifique.
    \item Rédigez un paragraphe de synthèse expliquant en quoi
          cet exercice illustre que le Trilemme CRO formalisé
          dans ce chapitre n'est pas une abstraction mathématique
          déconnectée, mais la \textbf{systématisation rigoureuse},
          désormais publiée et vérifiable, d'un raisonnement que
          vous avez pratiqué tout au long de ce cours.
\end{enumerate}
\end{boxexo}

\bigskip

\begin{boxinfo}
\textbf{Sur la place de ce chapitre dans le cours.} Ce chapitre
referme la boucle conceptuelle ouverte au tout premier chapitre
du cours~: le Trilemme CRO, présenté ici dans sa forme prouvée et
publiée, est l'outil que vous avez appliqué intuitivement à chaque
décision investigative rencontrée~--- du choix d'éteindre ou non
un PC suspect (Ch.~\ref{chap:custody}) à la stratégie de coopération
internationale (Ch.~\ref{chap:norm-glob}), en passant par la
conception d'une architecture de \textit{readiness} cloud-native
(Ch.~\ref{chap:containers}).

\textit{Que vous reteniez une seule chose de ce chapitre, que ce
soit celle-ci~: tout système, toute décision, toute primitive
implique un arbitrage entre confidentialité, fiabilité et
opposabilité~--- et ce n'est plus une intuition pédagogique mais
un théorème mathématique prouvé, publié et évalué par les pairs.
La compétence de l'investigateur ne consiste pas à éliminer cet
arbitrage~--- c'est désormais mathématiquement établi comme
structurellement impossible. Elle consiste à le rendre
\textbf{explicite, justifié et documenté}, et, quand l'institution
le permet, à investir dans $H(C)$ pour repousser la limite elle-même.}
\end{boxinfo}

% ============================================================================
% Références bibliographiques propres à ce chapitre
% (à intégrer dans D_bibliographie.tex)
% ============================================================================
% @article{minka2026clo,
%   author  = {Minka Mi Nguidjoï, Thierry Emmanuel and Mani Onana,
%              Flavien Serge and Djotio Ndjé, Thomas},
%   title   = {Cryptographic Legal Opposability (CLO) Framework:
%              Confidentiality-Reliability-Opposability Trade-Offs
%              in Institutional Cryptography},
%   journal = {IEEE Access},
%   volume  = {14},
%   pages   = {20427--20445},
%   year    = {2026},
%   doi     = {10.1109/ACCESS.2026.3660105}
% }
% @article{minka2026conditional,
%   author  = {Minka Mi Nguidjoï, Thierry Emmanuel and Mani Onana,
%              Flavien Serge and Djotio Ndjé, Thomas},
%   title   = {The Conditional CRO Trilemma: Impossibility Bounds
%              on Confidentiality, Reliability, and Opposability
%              Under Institutional Extraction Limits},
%   journal = {IEEE Access},
%   volume  = {14},
%   pages   = {33263--33274},
%   year    = {2026},
%   doi     = {10.1109/ACCESS.2026.3669065}
% }
% @article{shannon1948,
%   author  = {Shannon, Claude E.},
%   title   = {A Mathematical Theory of Communication},
%   journal = {Bell System Technical Journal},
%   volume  = {27}, number = {3}, pages = {379--423}, year = {1948}
% }
% @article{miller1956,
%   author  = {Miller, George A.},
%   title   = {The Magical Number Seven, Plus or Minus Two},
%   journal = {Psychological Review}, volume = {63}, pages = {81--97},
%   year    = {1956}
% }
% @inproceedings{adida2009helios,
%   author = {Adida, Ben}, title = {Helios: Web-based open-audit voting},
%   booktitle = {Proc. 17th USENIX Security Symp.}, pages = {335--348},
%   year = {2009}
% }
% @inproceedings{adida2009university,
%   author = {Adida, Ben and de Marneffe, Olivier and Pereira, Olivier
%             and Quisquater, Jean-Jacques},
%   title  = {Electing a university president using open-audit voting:
%             Analysis of real-world use of Helios},
%   booktitle = {Proc. EVT/WOTE}, year = {2009}
% }
% @misc{iacr2025,
%   author = {{International Association for Cryptologic Research}},
%   title  = {2025 IACR Election Update}, year = {2025},
%   note   = {Online; iacr.org/news/item/27138}
% }
% @article{johnson2001ecdsa,
%   author  = {Johnson, Don and Menezes, Alfred and Vanstone, Scott},
%   title   = {The Elliptic Curve Digital Signature Algorithm (ECDSA)},
%   journal = {Int. J. Inf. Secur.}, volume = {1}, number = {1},
%   pages   = {36--63}, year = {2001}
% }